\definition \textbf{Erwartungswert} (nicht-negativ)
$$
    \E[X] = \int_{0}^{\infty} \biggl( 1-F_X(x) \biggr)\ dx
$$
\subtext{$X: \Sigma \to \R,\quad X(\omega) \geq 0\ \forall \omega \in \Omega$}

{\scriptsize
    \remark $\E[X]$ kann unendlich sein
}

\theorem $\forall \omega \in \Omega: X(\omega) \geq 0 \implies \E[X] \geq 0$\\
\subtext{Gleichheit: $\E[X] = 0 \iff X=0$, fast sicher}

\definition \textbf{Erwartungswert} 
$$
    \E[X] = \E[X_+] - \E[X_-] \quad\text{falls}\quad \E[|X|]<\infty
$$
{\scriptsize
    \remark $X$ kein konst. Vorzeichen, nicht $\E[|X|] < \infty $: $\E[X]$ undefiniert
}

\subsection{Diskreter Erwartungswert}

\theorem \textbf{Diskreter Erwartungswert}
$$
    \E\bigl[ \phi(X) \bigr] = \sum_{w \in W} \phi(x) \cdot \P[X = x]
$$
\subtext{$X: \Omega \to \R,\quad W \cleq \N,\quad \phi: \R \to \R$}

\begin{center}
    \begin{tabular}{l|l|l}
        $\text{Ber}(p)$             & $\E[X] = p$                   & $\V[X] = p(1-p)$\\
        $\text{Poisson}(\lambda)$   & $\E[X] = \lambda$             & $\V[X] = \lambda$\\
        $\text{Bin}(n,p)$           & $\E[X] = n\cdot p$            & $\V[X] = np(1-p)$\\
        $\mathbb{I}_A$              & $\E[\mathbb{I}_A] = \P[A]$    \\
        $\exp(\lambda)$             & $\E[X] = \frac{1}{\lambda}$   \\
        $\mathcal{U}([a,b])$        & $\E[X] = \frac{a+b}{2}$ 
    \end{tabular}    
\end{center}

\newpage
\subsection{Stetiger Erwartungswert}

\definition \textbf{Erwartungswert} (stetig)
$$ 
    \E[X] = \int_{-\infty}^{\infty} x \cdot f(x)\ dx
$$
\subtext{$X: \Omega \to \R,\quad f(x) \text{ Dichtefunktion}$}

\theorem \textbf{Linearität}
\begin{align*}
    \text{(i)}  &\quad \E[\lambda X] &=& \lambda \E[X] \\
    \text{(ii)} &\quad \E[X + Y]     &=& \E[X] + \E[Y]
\end{align*}
\subtext{$X,Y:\Omega\to\R,\quad\lambda\in\R$}

\theorem \textbf{Monotonie}
$$
    X \leq Y \implies \E[X] \leq \E[Y]
$$

\theorem \textbf{Multiplikation}\\
\smalltext{$X,Y$ unabhängig}
$$
    \E[X \cdot Y] = \E[X]\cdot\E[Y]
$$

\theorem \textbf{Dichtefunktion bei Abbildungen}\\
\smalltext{$\phi:\R\to\R$ stückweise stetig, beschränkt} 
$$
    f \text{ Dichte von } X \iff \E[\phi(X)] = \int_{-\infty}^{\infty}\phi(x)f(x)\ dx
$$
\subtext{$f:\R\to\R_+$ s.d. $\int_{-\infty}^{\infty}f(x\ dx = 1)$}

\theorem \textbf{Unabhängigkeit} (durch Abbildungen)\\
\smalltext{$\forall \phi,\psi:\R\to\R$ stückweise stetig, beschränkt}
$$
    X,Y \text{ unabh.} \iff \E[\phi(X)\psi(Y)] = \E[\phi(X)]\cdot\E[\psi(Y)]
$$
\subtext{Auch verallgemeinert für $X_1,\ldots,X_n$, $\phi_1,\ldots,\phi_n$}

\newpage
\subsection{Ungleichungen}

\theorem \textbf{Markov}\\
\smalltext{$X \geq 0,\quad g: X(\Omega)\to[0,\infty)$ wachsend}
$$
    \forall c \in \R \text{ s.d. } g(c)>0:\qquad \P[X \geq c] \leq \frac{\E[g(X)]}{g(c)}
$$

\theorem \textbf{Jensen}\\
\smalltext{$\phi:\R\to\R$ konvex,$\quad \E[\phi(X)],\E[X]$ wohldefiniert}
$$
    \phi\Bigl(\E[X]\Bigr)\leq\E\Bigl[\phi(X)\Bigr]
$$

\lemma \textbf{Dreiecksungleichung}\\
\subtext{Jensen mit $\phi(X) = |X|$ und $\phi(X) = X^2$}
$$
    \Bigl\vert\E[X]\Bigr\vert\leq\E\Bigl[\vert X\vert\Bigr] \qquad \E[|X|] \leq \sqrt{\E[X^2]}
$$

\theorem \textbf{Chebychev}\\
\subtext{$Y$ s.d. $\V[Y] < \infty,\quad c>0$}
$$
    \P\Bigl[ |Y-\E[Y]| \geq c \Bigr] \leq \frac{\V[Y]}{c^2}
$$

\theorem \textbf{Chernoff}
% Slides, ende von ZGS

\newpage
\subsection{Varianz}

\definition \textbf{Varianz}\\ 
\subtext{$\E[X^2]<\infty$}
$$
    \mathbb{V}[X] := \E\Bigl[ (X - \E[X])^2 \Bigr]
$$
\definition \textbf{Standardabweichung}
$$
    \rho(X) := \sqrt{\mathbb{V}[X]}
$$
{\scriptsize
    \notation Auch $\rho, \rho_X$
}

\lemma \textbf{Varianz} (Alternativ)
$$
    \mathbb{V}[X] = \E[X^2]-\E[X]^2 
$$

{\footnotesize
    \textbf{Beispiel:} Bestimme $\E\bigl[X^2\bigr]$, wobei $X \sim \text{Poisson}(\lambda)$, $\lambda > 0$. 
    \begin{align*}
        \V[X]               &= \E\bigl[X^2\bigr] - \E[X]^2      & (X \sim \text{Poisson}(\lambda))      \\
        \lambda             &= \E\bigl[X^2\bigr] - \lambda^2                                            \\
        \lambda + \lambda^2 &= \E\bigl[X^2\bigr]                                                        \\
    \end{align*}
}

\lemma \textbf{Eigenschaften}
\begin{align*}
    \text{(i)}      &\quad \V[X]        \geq 0         \\
    \text{(ii)}     &\quad \V[aX]       = a^2\V[X]     \\
    \text{(iii)}    &\quad \V[X+a]      = \V[X]
\end{align*}

\lemma \textbf{Addition} (Unabhängigkeit)\\
\smalltext{$X_1,\ldots,X_n$ paarweise unabhängig}
$$
    \V\Biggl[ \sum_{k=1}^{n}X_k \Biggr] = \sum_{k=1}^{n}\V[X_k]
$$

\newpage
\subsection{Kovarianz}

\definition \textbf{Kovarianz}\\
\subtext{$X,Y$ s.d. $\E[X^2],\E[Y^2]<\infty$}
$$
    \text{cov}(X,Y) := \E\Bigl[ (X-\E[X])\cdot(Y-\E[Y]) \Bigr]
$$

\lemma \textbf{Kovarianz} (Alternativ)
$$
    \text{cov}(X,Y) = \E[XY] - \E[X]\E[Y]
$$

\remark $\text{cov}(X,X) = \V[X]$

\lemma \textbf{Eigenschaften von} $\text{cov}$
\begin{align*}
    \text{(i)}\quad         & \text{cov}(X,Y) \geq 0 \\
    \text{(ii)}\quad        & \text{cov}(X,Y) = \text{cov}(Y,X) \\
    \text{(iii)}\quad       & X,Y \text{ unabh. } \implies \text{cov}(X,Y) = 0\quad \text{\subtext{Nicht umgekehrt!}} \\
    \text{(iv)}\quad        & \V[X \pm Y] = \V[X] + \V[Y] \pm 2\text{cov}(X,Y) \\
    \text{(v)}\quad         & \text{cov}\Biggl( \sum_{i=1}^{n}X_i,\sum_{j=1}^{n}Y_i \Biggr) = \sum_{i=1}^{n}\sum_{j=1}^{n}\text{cov}(X_i,Y_j) \\
    \text{(vi)}\quad        & \text{cov}(aX+b, cY+d) = ac\cdot\text{cov}(X,Y) \\
    \text{(vii)}\quad       & \text{cov}\Bigl(X, (eY+f) + (gZ+h)\Bigr) = e\text{cov}(X,Y) + g\text{cov}(X,Z)
\end{align*}

\definition \textbf{Kovarianzmatrix} \subtext{$\quad \textbf{X} = (X_1,\ldots,X_n)^\top$}
{\small
    $$
    \Sigma = \text{cov}(\textbf{X}) = \begin{bmatrix}
        \V[X_1] & \text{cov}(X_1,X_2) & \cdots & \text{cov}(X_1,X_n) \\
        \text{cov}(X_2,X_1) & \text{cov}(X_2,X_2) & \cdots & \text{cov}(X_2,X_n) \\
        \vdots & \vdots & \ddots & \vdots \\
        \text{cov}(X_n,X_1) & \text{cov}(X_n,X_2) & \cdots & \text{cov}(X_n,X_n)
    \end{bmatrix}
    $$
}
